% ==========================================
% ANEXOS — ARCHIVO SEMI-ESTÁTICO
% ==========================================
% Agrega tus apéndices aquí. A diferencia de los otros
% archivos en common/, este SÍ se edita por proyecto.
%
% TIP: Los anexos se numeran con letras (A, B, C...).
%      Usa \chapter{} para cada anexo y \section{} para subdivisiones.
%
% TIP: Cada anexo que agregues aparecerá automáticamente
%      en el índice general.
% ==========================================

\appendix
\renewcommand{\appendixname}{Anexo}
\renewcommand{\appendixpagename}{Anexos}
\renewcommand{\appendixtocname}{Anexos}
\phantomsection
\addappheadtotoc
\appendixpage

\chapter{Códigos completos}
\label{chap:codigos_completos}
% Ejemplos:


A continuación se presentan los archivos que conforman el programa desarrollado para la comunicación con CUDA, incluyendo tanto la configuración de directivas en \texttt{Makefile}, como las definiciones estructurales de datos compartidos y la lógica central en lenguaje C. 

\section{Makefile}
\lstinputlisting[
	language=make,
	frame=single,
	basicstyle=\ttfamily\scriptsize,
	caption={Código de automatización de compilación y ejecución concurrente (Makefile). \propia{}},
	label={lst:codigo_makefile}
]{codigos_practica_8/Makefile}

\section{MultiplicacionMatrices.c}
\lstinputlisting[
	language=C,
	frame=single,
	basicstyle=\ttfamily\scriptsize,
	caption={Código de la multiplicación matricial serial en CPU (MultiplicacionMatrices.c). \propia{}},
	label={lst:codigo_c_serial}
]{codigos_practica_8/MultiplicacionMatrices.c}

\section{MultiplicacionMatricesCUDA.cu}
\lstinputlisting[
	language=C++,
	frame=single,
	basicstyle=\ttfamily\scriptsize,
	caption={Código de la multiplicación matricial paralela en GPU (MultiplicacionMatricesCUDA.cu). \propia{}},
	label={lst:codigo_cuda_paralelo}
]{codigos_practica_8/MultiplicacionMatricesCUDA.cu}

